\documentclass[12pt]{article}
\usepackage[utf8]{inputenc}
\usepackage[T1]{fontenc}
\usepackage[margin=1in]{geometry}
\usepackage{setspace}
\usepackage{xcolor}
\usepackage[colorlinks=true, urlcolor=blue!60!black, linkcolor=blue!60!black]{hyperref}
\setlength{\parindent}{0em}
\setlength{\parskip}{0.45em}
\renewcommand{\baselinestretch}{1.0}
\begin{document}
\begin{center}
{\Large \textit{Egbert v. Boule}, 596 U.S. 482 (2022)\par}
\vspace{0.3em}
{\normalsize Case Brief --- One-Page Edition\par}
\end{center}

\textbf{Memory Jogger.} No \textit{Bivens} suit against a Border Patrol agent who roughed up an innkeeper: if there is any reason to think Congress is better equipped to create a damages remedy, courts may not (\href{https://supreme.justia.com/cases/federal/us/596/21-147/}{Justia}).

\textbf{Facts.} Border Patrol Agent Egbert followed a guest into the driveway of Robert Boule's border-side inn and, when told to leave, allegedly threw Boule to the ground. After Boule complained, Egbert allegedly retaliated by prompting an IRS audit. Boule sued Egbert personally for Fourth Amendment excessive force and First Amendment retaliation.

\textbf{Procedural History.} The district court granted summary judgment for Egbert; the Ninth Circuit reversed.

\textbf{Votes.} Reversed: 6--3 (Fourth Amendment) and 9--0 (First Amendment). Thomas for the Court; Gorsuch concurred in the judgment; Sotomayor, with Breyer and Kagan, concurred in part and dissented in part.

\textbf{Holding.} \textit{Bivens} supplies no damages remedy for either claim. The two-step \textit{Bivens} test collapses into one question: whether there is any reason to think Congress might be better equipped to create a remedy.

\textbf{Reasoning.} Creating damages actions is legislative work. Border security and CBP's internal grievance process (even without damages or judicial review) each foreclosed the force claim; retaliation suits would chill legitimate conduct across the federal workforce. The Court declined to overrule \textit{Bivens}, which Gorsuch would have done. Sotomayor called the force claim ``materially indistinguishable'' from \textit{Bivens}. \emph{Later:} \textit{Goldey v. Fields} (2025) applied \textit{Egbert} to a prisoner's excessive-force claim.

\textbf{Examples.} \emph{Same-Side:} retaliation by a federal inspector is not actionable. \emph{Opposite-Side:} the same force by a state trooper is actionable under 42 U.S.C. \S~1983. \emph{Fence-Sitter:} a warrantless task-force arrest that closely tracks \textit{Bivens} itself.

\textbf{Critique.} State officers answer in damages for conduct federal officers now largely escape, and internal complaint processes can effectively immunize federal agents.

\textbf{Key Quotation.} ``While our cases describe two steps, those steps often resolve to a single question: whether there is any reason to think that Congress might be better equipped to create a damages remedy.'' --- Justice Thomas, 596 U.S. at 492.

\end{document}
